\documentclass[tikz,border=10pt]{standalone}
\usepackage{fontspec}
\usepackage[mathrm=sym]{unicode-math}
\setmainfont{texgyrepagella-regular.otf}[
    BoldFont = texgyrepagella-bold.otf,
    ItalicFont = texgyrepagella-italic.otf,
    BoldItalicFont = texgyrepagella-bolditalic.otf
]
\setmathfont{texgyrepagella-math.otf}[AutoFakeBold = 2.5]
\usepackage{tikz}
\usetikzlibrary{arrows.meta, calc}

\begin{document}
\begin{tikzpicture}[>=Stealth, line cap=round, line join=round]

  \def\w{9.6}
  \def\h{7.5}

  % Outer frame box
  \draw[line width=0.85pt] (0, 0) rectangle (\w, \h);

  % Bottom X-axis: rms-radius R (fm)
  \foreach \r in {2.2, 2.6, 3.0, 3.4} {
    \pgfmathsetmacro{\x}{(\r - 2.10)*6.0}
    \draw[line width=0.7pt] (\x, 0) -- (\x, 0.16);
    \draw[line width=0.7pt] (\x, \h) -- (\x, {\h - 0.16});
    \node[below=3pt, font=\small] at (\x, 0) {$\r$};
  }
  % Minor ticks at 2.0, 2.4, 2.8, 3.2, 3.6
  \foreach \r in {2.0, 2.4, 2.8, 3.2, 3.6} {
    \pgfmathsetmacro{\x}{(\r - 2.10)*6.0}
    \ifdim \x pt > 0 pt
      \ifdim \x pt < \w pt
        \draw[line width=0.5pt] (\x, 0) -- (\x, 0.10);
        \draw[line width=0.5pt] (\x, \h) -- (\x, {\h - 0.10});
      \fi
    \fi
  }
  \draw[->, line width=0.8pt] (2.2, -0.65) -- (6.0, -0.65) node[right=3pt, font=\normalsize] {$\text{rms-radius } R\ (\mathrm{fm})$};

  % Left Y-axis: Collective Mass M(R) (mass units)
  \foreach \m in {0, 5, 10, 15, 20, 25, 30, 35, 40, 45, 50} {
    \pgfmathsetmacro{\y}{\m * (7.5 / 50.0)}
    \draw[line width=0.7pt] (0, \y) -- (0.16, \y);
    \node[left=3pt, font=\small] at (0, \y) {$\m$};
  }
  \draw[->, line width=0.8pt] (-1.25, 0.5) -- (-1.25, 7.0);
  \node[rotate=90, font=\normalsize] at (-1.55, 3.75) {Collective Mass $M(R)\ (\text{mass units})$};

  % Right Y-axis: Collective Potential V(R) (MeV)
  \foreach \v in {-140, -120, -100, -80, -60} {
    \pgfmathsetmacro{\y}{(\v - (-140.0)) * (7.5 / 90.0)}
    \draw[line width=0.7pt] (\w, \y) -- ({\w - 0.16}, \y);
    \node[right=3pt, font=\small] at (\w, \y) {$\v$};
  }
  \draw[->, line width=0.8pt] ({\w + 1.3}, 0.5) -- ({\w + 1.3}, 7.0);
  \node[rotate=-90, font=\normalsize] at ({\w + 1.6}, 3.75) {Collective Potential $\mathcal{V}(R)\ (\mathrm{MeV})$};

  % Header Box: 16O SK-3 / ATDHF CHF
  \draw[line width=0.7pt] (2.2, 6.35) rectangle (4.4, 7.35);
  \node[font=\footnotesize\bfseries, align=center] at (3.3, 6.85) {$^{16}\mathrm{O}\quad \text{SK-3}$\\[2pt]$\text{ATDHF}\quad \text{CHF}$};

  % R_0 Arrows at R = 2.63 fm (x = (2.63 - 2.10)*6.0 = 3.18 cm)
  \def\xrzero{3.18}
  % Top arrow
  \node[font=\normalsize] at (\xrzero, 5.95) {$R_0$};
  \draw[->, line width=0.8pt] (\xrzero, 5.70) -- (\xrzero, 4.95);
  % Bottom arrow pointing down to axis from inside
  \draw[->, line width=0.8pt] (\xrzero, 0.75) -- (\xrzero, 0.05);

  % Horizontal levels: -107.7 MeV and -126.7 MeV
  \def\ylow{1.108}  % -126.7 MeV
  \def\yhigh{2.692} % -107.7 MeV
  \draw[line width=0.85pt] (1.3, \yhigh) -- (6.8, \yhigh) node[right=3pt, font=\small] {$-107.7$};
  \draw[line width=0.85pt] (1.5, \ylow) -- (6.8, \ylow) node[right=3pt, font=\small] {$-126.7$};

  % Clip curves to inside box
  \begin{scope}
    \clip (0, 0) rectangle (\w, \h);

    % 1. Collective Mass M(R) (solid)
    \draw[very thick, line width=1.1pt] plot[smooth] coordinates {
      (0.480, 7.350)
      (0.720, 6.450)
      (0.960, 5.625)
      (1.200, 5.025)
      (1.500, 4.425)
      (1.800, 3.975)
      (2.100, 3.675)
      (2.400, 3.480)
      (2.700, 3.420)
      (3.000, 3.450)
      (3.300, 3.570)
      (3.600, 3.780)
      (3.900, 4.020)
      (4.200, 4.320)
      (4.500, 4.575)
      (4.800, 4.740)
      (5.100, 4.800)
      (5.400, 4.830)
      (5.700, 4.875)
      (6.000, 4.980)
      (6.300, 5.100)
      (6.600, 5.250)
      (6.900, 5.430)
      (7.200, 5.670)
      (7.500, 5.925)
      (7.800, 6.225)
      (8.100, 6.570)
      (8.400, 6.930)
      (8.700, 7.320)
      (8.880, 7.500)
    };

    % 2. Collective Potential V(R) (dashed)
    \draw[thick, line width=1.0pt, dashed, dash pattern=on 4pt off 3pt] plot[smooth] coordinates {
      (0.720, 7.185)
      (0.960, 5.975)
      (1.200, 4.914)
      (1.440, 3.997)
      (1.680, 3.218)
      (1.920, 2.570)
      (2.160, 2.049)
      (2.400, 1.648)
      (2.640, 1.362)
      (2.880, 1.185)
      (3.120, 1.111)
      (3.360, 1.112)
      (3.600, 1.131)
      (3.840, 1.164)
      (4.080, 1.211)
      (4.320, 1.274)
      (4.560, 1.351)
      (4.800, 1.442)
      (5.040, 1.549)
      (5.280, 1.670)
      (5.520, 1.805)
      (5.760, 1.956)
      (6.000, 2.121)
      (6.240, 2.300)
      (6.480, 2.495)
      (6.720, 2.704)
      (6.960, 2.927)
      (7.200, 3.166)
      (7.440, 3.419)
      (7.680, 3.686)
      (7.920, 3.969)
      (8.160, 4.266)
      (8.400, 4.577)
      (8.640, 4.904)
      (8.880, 5.245)
      (9.120, 5.600)
      (9.360, 5.971)
    };

    % 3. Constrained HF Potential V_0(R) (dashed)
    \draw[thick, line width=1.0pt, dashed, dash pattern=on 4pt off 3pt] plot[smooth] coordinates {
      (1.080, 6.501)
      (1.320, 5.578)
      (1.560, 4.786)
      (1.800, 4.119)
      (2.040, 3.573)
      (2.280, 3.141)
      (2.520, 2.819)
      (2.760, 2.602)
      (3.000, 2.484)
      (3.240, 2.459)
      (3.480, 2.467)
      (3.720, 2.485)
      (3.960, 2.515)
      (4.200, 2.555)
      (4.440, 2.605)
      (4.680, 2.667)
      (4.920, 2.739)
      (5.160, 2.821)
      (5.400, 2.915)
      (5.640, 3.019)
      (5.880, 3.133)
      (6.120, 3.259)
      (6.360, 3.395)
      (6.600, 3.541)
      (6.840, 3.699)
      (7.080, 3.867)
      (7.320, 4.045)
      (7.560, 4.235)
      (7.800, 4.435)
      (8.040, 4.645)
      (8.280, 4.867)
      (8.520, 5.099)
      (8.760, 5.341)
      (9.000, 5.595)
      (9.240, 5.859)
    };
  \end{scope}

  % Labels and arrows to curves
  \draw[->, line width=0.7pt] (5.5, 4.8) node[right=2pt, font=\small] {$M(R)$} -- (4.7, 4.8);
  \draw[->, line width=0.7pt] (1.1, 1.4) node[left=1pt, font=\small] {$\mathcal{V}(R)$} -- (1.65, 1.35);
  \draw[->, line width=0.7pt] (6.0, 3.8) node[left=1pt, font=\small] {$\mathcal{V}_0(R)$} -- (6.6, 3.9);

  % Vertical double arrow for 19 MeV (drawn on top of curves with protective mask)
  \draw[<->, line width=0.8pt] (6.4, \ylow) -- (6.4, \yhigh) node[midway, left=2pt, font=\small\bfseries, fill=white, inner sep=1.2pt] {$19\ \mathrm{MeV}$};

\end{tikzpicture}
\end{document}
